
\documentclass{article}
\usepackage{colortbl}
\usepackage{makecell}
\usepackage{multirow}
\usepackage{supertabular}

\begin{document}

\newcounter{utterance}

\centering \large Interaction Transcript for game `function\_detective', experiment `string\_passive\_examples', episode 163 with Qwen3.6{-}35B{-}A3B{-}FP8{-}without{-}reasoning.
\vspace{24pt}

{ \footnotesize  \setcounter{utterance}{1}
\setlength{\tabcolsep}{0pt}
\begin{supertabular}{c@{$\;$}|p{.15\linewidth}@{}p{.15\linewidth}p{.15\linewidth}p{.15\linewidth}p{.15\linewidth}p{.15\linewidth}}
    \# & \multicolumn{2}{c}{Player} && \multicolumn{2}{c}{Game Master} \\
    \hline

    \theutterance \stepcounter{utterance}  
    & & & \multicolumn{4}{p{0.6\linewidth}}{
        \cellcolor[rgb]{0.9,0.9,0.9}{
            \makecell[{{p{\linewidth}}}]{
                \texttt{\tiny{[P1$\langle$GM]}}
                \texttt{I have a hidden Python function for you to figure out.} \\
\\ 
\texttt{Your job is to infer what it does and then submit a correct Python implementation.} \\
\\ 
\texttt{Function signature:} \\
\texttt{(text: str) {-}> str} \\
\\ 
\texttt{You have 15 turns total.} \\
\\ 
\texttt{I control which valid input{-}output examples you see. You can request them one at a time.} \\
\\ 
\texttt{You must make your best SOLVE attempt no later than the second{-}to{-}last turn.} \\
\texttt{Do not wait until the very end.} \\
\\ 
\texttt{You can use these actions:} \\
\texttt{{-} JSON format described at the end of the prompt} \\
\texttt{{-} SOLVE: ```python ... ```} \\
\\ 
\texttt{Do not add import statements in your guess.} \\
\texttt{Anything you might need is already available, so just write the function body in hypothesis and return the answer.} \\
\\ 
\texttt{Preloaded examples:} \\
\texttt{1. ('a{-}b', 'a{-}b')} \\
\\ 
\texttt{Scoring:} \\
\texttt{{-} correctness is required for any non{-}zero score} \\
\texttt{{-} in interactive modes, solving with fewer information requests is better} \\
\\ 
\texttt{Input formatting reminders:} \\
\texttt{{-} strings must be quoted} \\
\texttt{{-} booleans must be True/False} \\
\texttt{{-} use commas between multiple values} \\
\texttt{{-} do not write any extra action tags in the same message} \\
\\ 
\texttt{Format example:} \\
\texttt{You write:} \\
\texttt{\{'current\_mode': <probing/confirming/don't know; choose whichever best describes your current state in the task>, 'input\_rationale': <brief description of why you chose this hypothesis>, 'hypothesis': <either the exact string "unknown" if you have no guess yet OR your current best guess as runnable Python matching the given signature, with a one{-}line docstring stating what it does. Never write placeholder code, if you cannot commit to a guess, use "unknown". Match the given signature.} \\
\texttt{e.g.} \\
\texttt{def f(x):} \\
\texttt{	"""Returns the absolute value of x."""} \\
\texttt{	return x if x >= 0 else {-}x} \\
\\ 
\texttt{Use "unknown" if you have no guess.>, 'confidence\_score': <score between 1{-}100. 100 means highest confidence in your hypothesis>\}} \\
\\ 
\texttt{Output only one valid json per turn.} \\
            }
        }
    }
     \\ \\

    \theutterance \stepcounter{utterance}  
    & \multicolumn{4}{p{0.6\linewidth}}{
        \cellcolor[rgb]{0.561,0.737,0.561}{
            \makecell[{{p{\linewidth}}}]{
                \texttt{\tiny{[P1$\rangle$GM]}}
                \texttt{\{'current\_mode': 'probing', 'input\_rationale': 'I have only one example so far: ('a{-}b', 'a{-}b'). This suggests the function might be an identity function, or it might only modify strings under certain conditions (e.g., containing specific characters, being of a certain length, etc.). I need more examples to distinguish between identity and other possibilities.', 'hypothesis': 'unknown', 'confidence\_score': 10\}} \\
            }
        }
    }
    & & \\ \\

    \theutterance \stepcounter{utterance}  
    & & & \multicolumn{4}{p{0.6\linewidth}}{
        \cellcolor[rgb]{0.9,0.9,0.9}{
            \makecell[{{p{\linewidth}}}]{
                \texttt{\tiny{[P1$\langle$GM]}}
                \texttt{Game Violation: ParseError: Response is not valid JSON.} \\
\texttt{Output ONLY:} \\
\texttt{One JSON with confidence\_score, current\_mode, input\_rationale, and hypothesis keys ...} \\
\texttt{OR} \\
\texttt{SOLVE: ```python ... ```} \\
            }
        }
    }
     \\ \\

    \theutterance \stepcounter{utterance}  
    & \multicolumn{4}{p{0.6\linewidth}}{
        \cellcolor[rgb]{0.561,0.737,0.561}{
            \makecell[{{p{\linewidth}}}]{
                \texttt{\tiny{[P1$\rangle$GM]}}
                \texttt{\{"current\_mode": "probing", "input\_rationale": "I have only one example so far: ('a{-}b', 'a{-}b'). This suggests the function might be an identity function, or it might only modify strings under certain conditions (e.g., containing specific characters, being of a certain length, etc.). I need more examples to distinguish between identity and other possibilities.", "hypothesis": "unknown", "confidence\_score": 10\}} \\
            }
        }
    }
    & & \\ \\

    \theutterance \stepcounter{utterance}  
    & & & \multicolumn{4}{p{0.6\linewidth}}{
        \cellcolor[rgb]{0.9,0.9,0.9}{
            \makecell[{{p{\linewidth}}}]{
                \texttt{\tiny{[P1$\langle$GM]}}
                \texttt{OUTPUT: ('hello', 'llo')} \\
            }
        }
    }
     \\ \\

    \theutterance \stepcounter{utterance}  
    & \multicolumn{4}{p{0.6\linewidth}}{
        \cellcolor[rgb]{0.561,0.737,0.561}{
            \makecell[{{p{\linewidth}}}]{
                \texttt{\tiny{[P1$\rangle$GM]}}
                \texttt{\{"current\_mode": "probing", "input\_rationale": "The first example ('a{-}b', 'a{-}b') suggests identity for short strings or specific patterns. The second example ('hello', 'llo') shows that the first two characters are removed. This looks like a shift operation. Let's test if it always removes the first two characters, or if it depends on something else. I need an example with a string shorter than 2 characters or a string with different structure to confirm.", "hypothesis": "unknown", "confidence\_score": 20\}} \\
            }
        }
    }
    & & \\ \\

    \theutterance \stepcounter{utterance}  
    & & & \multicolumn{4}{p{0.6\linewidth}}{
        \cellcolor[rgb]{0.9,0.9,0.9}{
            \makecell[{{p{\linewidth}}}]{
                \texttt{\tiny{[P1$\langle$GM]}}
                \texttt{OUTPUT: ('   ', '   ')} \\
            }
        }
    }
     \\ \\

    \theutterance \stepcounter{utterance}  
    & \multicolumn{4}{p{0.6\linewidth}}{
        \cellcolor[rgb]{0.561,0.737,0.561}{
            \makecell[{{p{\linewidth}}}]{
                \texttt{\tiny{[P1$\rangle$GM]}}
                \texttt{\{"current\_mode": "probing", "input\_rationale": "Example 1: 'a{-}b' {-}> 'a{-}b' (length 3, unchanged). Example 2: 'hello' {-}> 'llo' (length 5, first 2 chars removed). Example 3: '   ' {-}> '   ' (length 3, unchanged). It seems like strings of length 3 are unchanged, but 'hello' (length 5) had its first 2 chars removed. Let's test a string of length 4 to see if it removes 2 chars or behaves differently. Also, let's test a string of length 1 or 2 to see edge cases.", "hypothesis": "unknown", "confidence\_score": 30\}} \\
            }
        }
    }
    & & \\ \\

    \theutterance \stepcounter{utterance}  
    & & & \multicolumn{4}{p{0.6\linewidth}}{
        \cellcolor[rgb]{0.9,0.9,0.9}{
            \makecell[{{p{\linewidth}}}]{
                \texttt{\tiny{[P1$\langle$GM]}}
                \texttt{OUTPUT: ('aba', 'aba')} \\
            }
        }
    }
     \\ \\

    \theutterance \stepcounter{utterance}  
    & \multicolumn{4}{p{0.6\linewidth}}{
        \cellcolor[rgb]{0.561,0.737,0.561}{
            \makecell[{{p{\linewidth}}}]{
                \texttt{\tiny{[P1$\rangle$GM]}}
                \texttt{\{"current\_mode": "probing", "input\_rationale": "Examples so far: 'a{-}b'{-}>'a{-}b', 'hello'{-}>'llo', '   '{-}>'   ', 'aba'{-}>'aba'. It seems strings of length 3 are unchanged. 'hello' (length 5) had first 2 chars removed. Let's test a string of length 4 to see if it removes 2 chars (resulting in length 2) or does something else. Also, let's test a string of length 1 or 2 to see if they are unchanged or modified.", "hypothesis": "unknown", "confidence\_score": 40\}} \\
            }
        }
    }
    & & \\ \\

    \theutterance \stepcounter{utterance}  
    & & & \multicolumn{4}{p{0.6\linewidth}}{
        \cellcolor[rgb]{0.9,0.9,0.9}{
            \makecell[{{p{\linewidth}}}]{
                \texttt{\tiny{[P1$\langle$GM]}}
                \texttt{OUTPUT: ('ab', 'ab')} \\
            }
        }
    }
     \\ \\

    \theutterance \stepcounter{utterance}  
    & \multicolumn{4}{p{0.6\linewidth}}{
        \cellcolor[rgb]{0.561,0.737,0.561}{
            \makecell[{{p{\linewidth}}}]{
                \texttt{\tiny{[P1$\rangle$GM]}}
                \texttt{\{"current\_mode": "probing", "input\_rationale": "Examples so far: 'a{-}b'{-}>'a{-}b' (len 3), 'hello'{-}>'llo' (len 5, first 2 removed), '   '{-}>'   ' (len 3), 'aba'{-}>'aba' (len 3), 'ab'{-}>'ab' (len 2). It seems strings of length <= 3 are unchanged. 'hello' (len 5) had first 2 chars removed. Let's test a string of length 4 to see if it removes 2 chars (resulting in length 2) or does something else. Also, let's test a longer string like 'abcdef' to see if it removes 2 chars.", "hypothesis": "unknown", "confidence\_score": 50\}} \\
            }
        }
    }
    & & \\ \\

    \theutterance \stepcounter{utterance}  
    & & & \multicolumn{4}{p{0.6\linewidth}}{
        \cellcolor[rgb]{0.9,0.9,0.9}{
            \makecell[{{p{\linewidth}}}]{
                \texttt{\tiny{[P1$\langle$GM]}}
                \texttt{OUTPUT: (' ', ' ')} \\
            }
        }
    }
     \\ \\

    \theutterance \stepcounter{utterance}  
    & \multicolumn{4}{p{0.6\linewidth}}{
        \cellcolor[rgb]{0.561,0.737,0.561}{
            \makecell[{{p{\linewidth}}}]{
                \texttt{\tiny{[P1$\rangle$GM]}}
                \texttt{\{"current\_mode": "probing", "input\_rationale": "Examples so far: 'a{-}b'{-}>'a{-}b' (len 3), 'hello'{-}>'llo' (len 5, first 2 removed), '   '{-}>'   ' (len 3), 'aba'{-}>'aba' (len 3), 'ab'{-}>'ab' (len 2), ' '{-}>' ' (len 1). It seems strings of length <= 3 are unchanged. 'hello' (len 5) had first 2 chars removed. Let's test a string of length 4 to see if it removes 2 chars (resulting in length 2) or does something else. Also, let's test a longer string like 'abcdef' to see if it removes 2 chars.", "hypothesis": "unknown", "confidence\_score": 60\}} \\
            }
        }
    }
    & & \\ \\

    \theutterance \stepcounter{utterance}  
    & & & \multicolumn{4}{p{0.6\linewidth}}{
        \cellcolor[rgb]{0.9,0.9,0.9}{
            \makecell[{{p{\linewidth}}}]{
                \texttt{\tiny{[P1$\langle$GM]}}
                \texttt{OUTPUT: ('hello world', 'rld')} \\
            }
        }
    }
     \\ \\

    \theutterance \stepcounter{utterance}  
    & \multicolumn{4}{p{0.6\linewidth}}{
        \cellcolor[rgb]{0.561,0.737,0.561}{
            \makecell[{{p{\linewidth}}}]{
                \texttt{\tiny{[P1$\rangle$GM]}}
                \texttt{\{"current\_mode": "probing", "input\_rationale": "Let's analyze the pattern: 'hello' (len 5) {-}> 'llo' (removed first 2). 'hello world' (len 11) {-}> 'rld' (removed first 9). This doesn't look like a simple fixed removal. Let's look at the indices. 'hello': indices 0,1 removed. 'hello world': indices 0{-}8 removed. The number of removed characters seems to be related to the length. For len 5, removed 2. For len 11, removed 9. Let's check if it's removing the first (len{-}3) characters? For len 5: 5{-}3=2. For len 11: 11{-}3=8. But 'rld' is 3 chars, so 11{-}3=8 removed, leaving 3. Wait, 'hello world' is 11 chars. 'rld' is 3 chars. 11{-}3=8. So 8 chars removed. For 'hello', 5{-}3=2 chars removed. For len 3, 3{-}3=0 chars removed. For len 2, 2{-}3={-}1? No, 0 removed. For len 1, 0 removed. For len 0, 0 removed. So it seems to keep the} \\
            }
        }
    }
    & & \\ \\

    \theutterance \stepcounter{utterance}  
    & & & \multicolumn{4}{p{0.6\linewidth}}{
        \cellcolor[rgb]{0.9,0.9,0.9}{
            \makecell[{{p{\linewidth}}}]{
                \texttt{\tiny{[P1$\langle$GM]}}
                \texttt{Game Violation: ParseError: Response is not valid JSON.} \\
\texttt{Output ONLY:} \\
\texttt{One JSON with confidence\_score, current\_mode, input\_rationale, and hypothesis keys ...} \\
\texttt{OR} \\
\texttt{SOLVE: ```python ... ```} \\
            }
        }
    }
     \\ \\

    \theutterance \stepcounter{utterance}  
    & \multicolumn{4}{p{0.6\linewidth}}{
        \cellcolor[rgb]{0.561,0.737,0.561}{
            \makecell[{{p{\linewidth}}}]{
                \texttt{\tiny{[P1$\rangle$GM]}}
                \texttt{\{"current\_mode": "probing", "input\_rationale": "Based on examples: 'a{-}b' (len 3) {-}> 'a{-}b' (keep last 3), 'hello' (len 5) {-}> 'llo' (keep last 3), '   ' (len 3) {-}> '   ' (keep last 3), 'aba' (len 3) {-}> 'aba' (keep last 3), 'ab' (len 2) {-}> 'ab' (keep last 2, since len < 3), ' ' (len 1) {-}> ' ' (keep last 1, since len < 3), 'hello world' (len 11) {-}> 'rld' (keep last 3). The pattern is: return the last 3 characters of the string, or the entire string if its length is less than 3.", "hypothesis": "unknown", "confidence\_score": 90\}} \\
            }
        }
    }
    & & \\ \\

    \theutterance \stepcounter{utterance}  
    & & & \multicolumn{4}{p{0.6\linewidth}}{
        \cellcolor[rgb]{0.9,0.9,0.9}{
            \makecell[{{p{\linewidth}}}]{
                \texttt{\tiny{[P1$\langle$GM]}}
                \texttt{OUTPUT: ('A', 'A')} \\
            }
        }
    }
     \\ \\

    \theutterance \stepcounter{utterance}  
    & \multicolumn{4}{p{0.6\linewidth}}{
        \cellcolor[rgb]{0.561,0.737,0.561}{
            \makecell[{{p{\linewidth}}}]{
                \texttt{\tiny{[P1$\rangle$GM]}}
                \texttt{\{"current\_mode": "confirming", "input\_rationale": "The hypothesis is that the function returns the last 3 characters of the string, or the entire string if its length is less than 3. The new example ('A', 'A') confirms this: length is 1, so it returns the whole string. All previous examples are consistent with this rule. I am confident enough to solve.", "hypothesis": "def f(text: str) {-}> str:\textbackslash{}n    \textbackslash{}"\textbackslash{}"\textbackslash{}"Returns the last 3 characters of the string, or the entire string if its length is less than 3.\textbackslash{}"\textbackslash{}"\textbackslash{}"\textbackslash{}n    if len(text) < 3:\textbackslash{}n        return text\textbackslash{}n    return text{[}{-}3:{]}", "confidence\_score": 95\}} \\
            }
        }
    }
    & & \\ \\

    \theutterance \stepcounter{utterance}  
    & & & \multicolumn{4}{p{0.6\linewidth}}{
        \cellcolor[rgb]{0.9,0.9,0.9}{
            \makecell[{{p{\linewidth}}}]{
                \texttt{\tiny{[P1$\langle$GM]}}
                \texttt{OUTPUT: ('aA', 'aA')} \\
            }
        }
    }
     \\ \\

    \theutterance \stepcounter{utterance}  
    & \multicolumn{4}{p{0.6\linewidth}}{
        \cellcolor[rgb]{0.561,0.737,0.561}{
            \makecell[{{p{\linewidth}}}]{
                \texttt{\tiny{[P1$\rangle$GM]}}
                \texttt{\{"current\_mode": "confirming", "input\_rationale": "The new example ('aA', 'aA') is consistent with the hypothesis: length is 2, so it returns the whole string. All examples so far support the rule: return the last 3 characters if length >= 3, otherwise return the entire string.", "hypothesis": "def f(text: str) {-}> str:\textbackslash{}n    \textbackslash{}"\textbackslash{}"\textbackslash{}"Returns the last 3 characters of the string, or the entire string if its length is less than 3.\textbackslash{}"\textbackslash{}"\textbackslash{}"\textbackslash{}n    if len(text) < 3:\textbackslash{}n        return text\textbackslash{}n    return text{[}{-}3:{]}", "confidence\_score": 98\}} \\
            }
        }
    }
    & & \\ \\

    \theutterance \stepcounter{utterance}  
    & & & \multicolumn{4}{p{0.6\linewidth}}{
        \cellcolor[rgb]{0.9,0.9,0.9}{
            \makecell[{{p{\linewidth}}}]{
                \texttt{\tiny{[P1$\langle$GM]}}
                \texttt{OUTPUT: ('a', 'a')} \\
            }
        }
    }
     \\ \\

    \theutterance \stepcounter{utterance}  
    & \multicolumn{4}{p{0.6\linewidth}}{
        \cellcolor[rgb]{0.561,0.737,0.561}{
            \makecell[{{p{\linewidth}}}]{
                \texttt{\tiny{[P1$\rangle$GM]}}
                \texttt{\{"current\_mode": "confirming", "input\_rationale": "The new example ('a', 'a') is consistent with the hypothesis: length is 1, so it returns the whole string. All examples so far support the rule: return the last 3 characters if length >= 3, otherwise return the entire string. I am confident enough to solve.", "hypothesis": "def f(text: str) {-}> str:\textbackslash{}n    \textbackslash{}"\textbackslash{}"\textbackslash{}"Returns the last 3 characters of the string, or the entire string if its length is less than 3.\textbackslash{}"\textbackslash{}"\textbackslash{}"\textbackslash{}n    if len(text) < 3:\textbackslash{}n        return text\textbackslash{}n    return text{[}{-}3:{]}", "confidence\_score": 99\}} \\
            }
        }
    }
    & & \\ \\

    \theutterance \stepcounter{utterance}  
    & & & \multicolumn{4}{p{0.6\linewidth}}{
        \cellcolor[rgb]{0.9,0.9,0.9}{
            \makecell[{{p{\linewidth}}}]{
                \texttt{\tiny{[P1$\langle$GM]}}
                \texttt{OUTPUT: ('123', '123')} \\
            }
        }
    }
     \\ \\

    \theutterance \stepcounter{utterance}  
    & \multicolumn{4}{p{0.6\linewidth}}{
        \cellcolor[rgb]{0.561,0.737,0.561}{
            \makecell[{{p{\linewidth}}}]{
                \texttt{\tiny{[P1$\rangle$GM]}}
                \texttt{\{"current\_mode": "confirming", "input\_rationale": "The new example ('123', '123') is consistent with the hypothesis: length is 3, so it returns the last 3 characters (which is the whole string). All examples so far support the rule: return the last 3 characters if length >= 3, otherwise return the entire string. I am confident enough to solve.", "hypothesis": "def f(text: str) {-}> str:\textbackslash{}n    \textbackslash{}"\textbackslash{}"\textbackslash{}"Returns the last 3 characters of the string, or the entire string if its length is less than 3.\textbackslash{}"\textbackslash{}"\textbackslash{}"\textbackslash{}n    if len(text) < 3:\textbackslash{}n        return text\textbackslash{}n    return text{[}{-}3:{]}", "confidence\_score": 100\}} \\
            }
        }
    }
    & & \\ \\

    \theutterance \stepcounter{utterance}  
    & & & \multicolumn{4}{p{0.6\linewidth}}{
        \cellcolor[rgb]{0.9,0.9,0.9}{
            \makecell[{{p{\linewidth}}}]{
                \texttt{\tiny{[P1$\langle$GM]}}
                \texttt{OUTPUT: ('', '')} \\
            }
        }
    }
     \\ \\

    \theutterance \stepcounter{utterance}  
    & \multicolumn{4}{p{0.6\linewidth}}{
        \cellcolor[rgb]{0.561,0.737,0.561}{
            \makecell[{{p{\linewidth}}}]{
                \texttt{\tiny{[P1$\rangle$GM]}}
                \texttt{\{"current\_mode": "confirming", "input\_rationale": "The new example ('', '') is consistent with the hypothesis: length is 0, so it returns the whole string. All examples so far support the rule: return the last 3 characters if length >= 3, otherwise return the entire string. I am confident enough to solve.", "hypothesis": "def f(text: str) {-}> str:\textbackslash{}n    \textbackslash{}"\textbackslash{}"\textbackslash{}"Returns the last 3 characters of the string, or the entire string if its length is less than 3.\textbackslash{}"\textbackslash{}"\textbackslash{}"\textbackslash{}n    if len(text) < 3:\textbackslash{}n        return text\textbackslash{}n    return text{[}{-}3:{]}", "confidence\_score": 100\}} \\
            }
        }
    }
    & & \\ \\

    \theutterance \stepcounter{utterance}  
    & & & \multicolumn{4}{p{0.6\linewidth}}{
        \cellcolor[rgb]{0.9,0.9,0.9}{
            \makecell[{{p{\linewidth}}}]{
                \texttt{\tiny{[P1$\langle$GM]}}
                \texttt{OUTPUT: ('bXROc35M9C3F', 'C3F')} \\
            }
        }
    }
     \\ \\

    \theutterance \stepcounter{utterance}  
    & \multicolumn{4}{p{0.6\linewidth}}{
        \cellcolor[rgb]{0.561,0.737,0.561}{
            \makecell[{{p{\linewidth}}}]{
                \texttt{\tiny{[P1$\rangle$GM]}}
                \texttt{\{"current\_mode": "confirming", "input\_rationale": "The new example ('bXROc35M9C3F', 'C3F') is consistent with the hypothesis: length is 12, so it returns the last 3 characters ('C3F'). All examples so far support the rule: return the last 3 characters if length >= 3, otherwise return the entire string. I am confident enough to solve.", "hypothesis": "def f(text: str) {-}> str:\textbackslash{}n    \textbackslash{}"\textbackslash{}"\textbackslash{}"Returns the last 3 characters of the string, or the entire string if its length is less than 3.\textbackslash{}"\textbackslash{}"\textbackslash{}"\textbackslash{}n    if len(text) < 3:\textbackslash{}n        return text\textbackslash{}n    return text{[}{-}3:{]}", "confidence\_score": 100\}} \\
            }
        }
    }
    & & \\ \\

    \theutterance \stepcounter{utterance}  
    & & & \multicolumn{2}{p{0.3\linewidth}}{
        \cellcolor[rgb]{0.9,0.9,0.9}{
            \makecell[{{p{\linewidth}}}]{
                \texttt{\tiny{[GM$|$GM]}}
                \texttt{Game Verdict: LOSS (turn limit reached)} \\
            }
        }
    }
    & & \\ \\

\end{supertabular}
}

\end{document}
